\section{Introdução}
\label{sec:introducao}

O crescimento do volume de dados tem impulsionado a adoção de arquiteturas escaláveis para análise e tomada de decisão \cite{ibmbigdata}. Nesse contexto, o desafio não está apenas em armazenar grandes volumes de dados de forma eficiente, mas também em processá-los equilibrando custos e desempenho \cite{weintraub2024}. A arquitetura de \textit{data lakes} baseada em armazenamento de objetos oferece escalabilidade e custos reduzidos, sendo que a separação entre camadas de armazenamento e processamento permite utilização de recursos computacionais em modo \textit{on-demand} \cite{dremiodatalake,weintraub2024}.

Ferramentas como o \textit{Apache Iceberg} têm se destacado ao permitir a gestão eficiente de tabelas em \textit{data lakes}, oferecendo versionamento, \textit{schema evolution} e melhorias de consulta \cite{iceberg2024}. Mecanismos de execução distribuída, como o \textit{Trino}, viabilizam consultas de alta \textit{performance} sobre dados armazenados em sistemas de objetos, por meio de processamento distribuído entre múltiplos nós trabalhadores \cite{trinoconcepts}. No entanto, a escolha do ambiente de execução das consultas ainda tende a ser estática, sendo que sistemas tradicionais não adaptam automaticamente a infraestrutura de processamento às características específicas de cada consulta \cite{weintraub2024,xue2024}. Trabalhos recentes têm explorado roteamento baseado em modelos de custo aprendidos \cite{strausz2025} e elasticidade em tempo de execução para análise de dados em nuvem \cite{zhang2025}.

A pesquisa propõe o \textit{DyraSQL}, uma abordagem adaptativa para direcionamento de consultas SQL, considerando características previamente disponíveis nos metadados. A proposta consiste em analisar a complexidade das consultas e o volume efetivo de dados antes da execução e, a partir disso, direcioná-las para \textit{clusters} \textit{Trino} de capacidades distintas, classificados como \textit{small}, \textit{medium} e \textit{large}. Tal abordagem busca viabilizar uma utilização mais eficiente dos recursos disponíveis, equilibrando custo e \textit{performance}. Xue et al. (2024) demonstram a importância de execução adaptativa de consultas em ambientes de \textit{data lakehouse} \cite{xue2024}, reforçando a relevância desta pesquisa.

A capacidade de direcionar automaticamente consultas para o ambiente mais adequado resulta em melhor utilização de recursos computacionais e em redução de provisionamento excessivo para consultas leves. Este trabalho desenvolve o \textit{DyraSQL}, um mecanismo de decisão que utiliza informações de metadados para definir, de forma automática, o \textit{cluster} mais adequado para a execução de consultas SQL. Para alcançar este propósito, o trabalho mapeia os metadados disponíveis em tabelas \textit{Apache Iceberg}, identificando quais podem ser utilizados como indicadores de volume de dados, complexidade da consulta e custo esperado de execução. Estabelece regras de decisão para roteamento, definindo parâmetros que permitam classificar consultas nas três faixas de capacidade. Projeta e desenvolve um protótipo funcional do \textit{DyraSQL} em ambiente local, utilizando o \textit{Trino} como \textit{engine} de consultas, Docker Compose para orquestração, MinIO para armazenamento de objetos e PostgreSQL para \textit{cache} e histórico. Realiza experimentos controlados nesse ambiente, a fim de verificar a viabilidade técnica do \textit{DyraSQL} e coletar métricas de consistência dos \textit{scores}, latência do \textit{EXPLAIN (TYPE IO)} e comportamento do \textit{cache}. Compara o comportamento do roteamento dinâmico com a lógica estática de um único \textit{cluster}, avaliando a capacidade do algoritmo de diferenciar consultas leves, intermediárias e pesadas.

Este trabalho está organizado da forma apresentada. A Seção~\ref{sec:referencial-teorico} aborda o referencial teórico. A Seção~\ref{sec:revisao-bibliografica} apresenta a revisão bibliográfica. A Seção~\ref{sec:metodologia} detalha a metodologia. A Seção~\ref{sec:experimentos} apresenta os experimentos e resultados. A Seção~\ref{sec:consideracoes-finais} apresenta as considerações finais e trabalhos futuros.
